\subsection{有限生成代数的维数}

命题 3. --- 设 \(\rho: A \to B\) 是一个环同态。置 \(n = \sup_{\mathfrak{p} \in \operatorname{Spec}(A)} \dim(B \otimes_A \kappa(\mathfrak{p}))\)。我们有不等式

\[
\dim (\mathbf {B}) + 1 \leqslant (\dim (\mathbf {A}) + 1). (n + 1).
\]

我们可以假设 \(\dim(A) \neq -\infty\) 且 \(n < +\infty\)。设 \(\mathfrak{q}_0 \subset \ldots \subset \mathfrak{q}_m\) 是 B 的素理想链；置 \(\mathfrak{p}_i = \rho^{-1}(\mathfrak{q}_i)\)。序列 \(\mathfrak{p}_i\) 是递增的，因此其值的集合的基数 \(\leqslant \dim(A) + 1\)。对于每个 \(\mathfrak{p} \in \operatorname{Spec}(A)\)，满足 \(\mathfrak{q}_j\) 且 \(\mathfrak{p}_j = \mathfrak{p}\) 的项构成 \(({}^{a}\rho)^{-1}(\mathfrak{p})\) 的子集 \(\operatorname{Spec}(B)\) 中的一条链，因而其基数小于等于 \(\dim(B \otimes_A \kappa(\mathfrak{p})) + 1\)（no. 1, 引理 1），并且因而小于等于 \(n^o 1\)\(n + 1\)。由此 \(m + 1 \leqslant (\dim(A) + 1)(n + 1)\)，从而得到命题。

注 1. --- 如果环 A 和 B 是诺特环，我们将在下文（§ 3, no. 4, 命题 7 的推论 2）看到，有不等式 dim(B) ≤ dim(A) + n，它强于命题 3 的不等式。

推论 1. --- 假设有  ，并且存在一个整数 n，使得对于所有 \textless{}，\_ 满足 \textless{}。则有

推论 2. --- 设 A 是一个环，{}{} 是系数在 A 中的一个不定元多项式环。我们有：

\[
1 + \dim (\mathrm{A}) \leqslant \dim (\mathrm{B}) \leqslant 1 + 2 \dim (\mathrm{A}).
\]

第一个不等式已经得到证明（§ 1, no. 3, 例 4）。现在证明第二个。对于 A 的每个素理想  ，环 {（它同构于 }{}）是主环且不是域，因而维数为 1（§ 1, no. 3, 例 2），不等式由命题 3 得出。

注 2. --- 我们将在稍后（§ 3, no. 4, 命题 7 的推论 3）看到，如果 A 是诺特环，则有 {}{}。

然而，无论取整数 \(n\) 和 \(q\)，只要满足 \(n + 1 \leqslant q \leqslant 2n + 1\)，总存在一个维数为 \(n\) 的环 A，使得 \(\dim(A[X]) = q\)（见 p. 84, 练习 7）。

推论 3. --- 如果 A 是有限维的，则每个非零有限型 A-代数都是有限维的。

事实上，由推论 2 对 \(n\) 作归纳可知，环 \(\mathbf{A}[\mathbf{T}_{1}, \ldots, \mathbf{T}_{n}]\) 是有限维的；若 \(\mathbf{A}\) 是有限维的，则特别地，\(\mathbf{A}[\mathbf{T}_{1}, \ldots, \mathbf{T}_{n}]\) 的每个非零商环\(\S\) 都是有限维的（§ 1, no. 3, 命题 6）。

